% !PRINTED 291
\weberpage{291}

To these invariants belongs also the degree $n$ of the group $P$ itself, which is
determined by the $j$ as follows:
\weq{2}{n = j_1 j_2 \cdots j_{\mu-1} j_\mu.}
For, by the meaning of the indices, $n : j_1$ is the degree of $P_1$, $n : j_1 j_2$
the degree of $P_2$, and so on, and since the last of the degrees is equal to $1$,
formula (2) results.

In general, in a composition series of $P$ each term is a normal divisor only of
the one immediately preceding it. But there may also occur individual terms which
are normal divisors of terms preceding them further back as well. Especially
important are those terms which are normal divisors of $P$ itself, and hence also
of all the terms of the composition series that precede them. We shall see what
algebraic consequences are to be drawn from this circumstance.

Let $P_\nu$ be a term of the series (1) which is at the same time a normal divisor
of $P$. The index of $P_\nu$ with respect to $P$ is $j_1 j_2 \ldots j_\nu$, and
this product is the degree of the partial resolvent by which the group $P$ is
reduced to $P_\nu$, which we shall denote by $\chi(y) = 0$. The group of this
resolvent, which is a normal equation, we obtain by Theorem 1. in the form
$P/P_\nu$.

We can now easily find a composition series for this group, together with the
corresponding index series, namely:
\weq{3}{
\begin{gathered}
P/P_\nu,\ P_1/P_\nu,\ P_2/P_\nu,\ \ldots,\ P_{\nu-1}/P_\nu,\ 1 \\
j_1,\quad j_2,\quad \ldots,\quad j_{\nu-1},\quad j_\nu.
\end{gathered}
}
For, by Theorem 1., \S.\,6, $P_1/P_\nu$ is a normal divisor of $P/P_\nu$ of
index $j_1$, and it is a greatest normal divisor,
because by 2., \S.\,6 no normal divisor of $P/P_\nu$ can stand above $P_1/P_\nu$
if no normal divisor of $P$ stands above $P_1$. And one can conclude in the same
way with respect to the following terms of (3).

From this we draw a further important consequence. We have already shown in
\S.\,7 that, if $Q$ is any normal divisor of $P$, a composition series of $P$ can
be found in which $Q$ occurs. If now $R$ is another normal divisor of $P$ and at
the same time a divisor of $Q$, then the composition series of $P$ can be so
arranged that $Q$ and $R$ occur in it.

We suppose such a composition series to be determined:
\weq{4}{P,\ Q',\ Q'',\ \ldots,\ Q,\ \ldots,\ R.}
